\subsection{子模的不变因子}

定理1. --- 设L是主理想整环A上的自由模，且设M是L的秩 n 为有限的子模。则存在L的一组基B、B 中的 n 个元素 \(e_{i}\)，以及 A 的 n 个非零元素 \(\alpha_{i}\)（\(1 \leqslant i \leqslant n\)），使得：

\begin{enumerate}
\def\labelenumi{\alph{enumi})}
\item
  诸 \(\alpha_{i}e_{i}\) 构成 \(\mathbf{M}\) 的一组基;
\item
  \(\alpha_{i}\) 整除 \(\alpha_{i + 1}\)，其中 \(1\leqslant i\leqslant n - 1\)
\end{enumerate}

此外，由 \((e_{i})\) 生成的模 M' 以及诸主理想 \(A\alpha_{i}\) 都由上述条件唯一确定；模 M'/M 是 L/M 的挠子模，并且同构于诸 A-模 A/A \(\alpha_{i}\) 的直和；最后，L/M 是 M'/M 与一个同构于 L/M' 的自由模的直和。

\begin{enumerate}
\def\labelenumi{\arabic{enumi})}
\tightlist
\item
  \(e_i\) 与 \(\alpha_i\) 的存在性.
\end{enumerate}

如果 \(M = \{0\}\)，那么定理是平凡的。如果 \(M \neq \{0\}\)，则由命题3可知，存在 L 的一个元素 \(e_{1}\)、A 的一个非零元素 \(\alpha_{1}\) 以及 L 的一个子模 \(L_{1}\)，使得 L 是 \(Ae_{1}\) 与 \(L_{1}\) 的直和，使得 M 是 \(A\alpha_{1}e_{1}\) 与 L 的子模 \(M_{1} = M \cap L_{1}\) 的直和，并且使得对 L 上的每个线性形式 g 都有 \(g(M) \subset A\alpha_{1}\)。

现在我们可以对 \(n\)（即 \(M\) 的秩）进行归纳。由于 \(L_1\) 是自由模（VII，第 15 页，推论 2），且 \(M_1\) 的秩为 \(n - 1\)，故存在一组基 \(B_1\)（\(L_1\) 的）、\((n - 1)\) 个元素 \(e_2, \ldots, e_n\)（取自 \(B_1\)），以及非零元素 \(\alpha_2, \ldots, \alpha_n\)（\(A\) 中的），使得 \((\alpha_2e_2, \ldots, \alpha_ne_n)\) 是 \(M_1\) 的一组基，且 \(\alpha_i\) 整除 \(\alpha_{i+1}\)，其中 \(2 \leqslant i \leqslant n - 1\)。若 \(L'\) 是 \(L_1\) 的由 \(B_1\) 中异于 \(e_2, \ldots, e_n\) 的元素生成的子模，则 \(L\) 是 \(L'\) 与模 \(M'\)（由 \(e_1, \ldots, e_n\) 生成的）的直和；而 \((e_1, \ldots, e_n)\) 是 \(M'\) 的一组基，且 \((\alpha_1e_1, \ldots, \alpha_ne_n)\) 是 \(M\) 的一组基。现在只需证明 \(\alpha_1\) 整除 \(\alpha_2\)；但 \(A\alpha_2\) 具有形式 \(g(M_1)\)，其中 \(g\) 是 \(L\) 上由 \(g(e_2) = 1\)、\(g(e_i) = 0\)（\(i \neq 2\)）以及 \(g(L') = \{0\}\) 定义的线性形式；而我们上面已经看到 \(g(M_1) \subset A\alpha_1\).

2) 唯一性性质。

由于诸 \(\alpha_{i}\) 非零，模 \(\mathbf{M}'\) 是由那些 \(x\in L\)（使得 \(\beta x\in M\) 对 A 中某个 \(\beta \neq 0\) 成立）组成的集合；换言之，\(\mathbf{M}' / \mathbf{M}\) 是 L/M 的挠子模。这就唯一确定了 \(\mathbf{M}'\)。

显然 \(M^{\prime}/M\) 同构于 n 个循环模 \(A/A\alpha_{i}\) 的直和（II，第 204 页，公式 (26)）。设 r 为异于 A 的理想 \(A\alpha_{i}\) 的个数：于是前 n−r 个理想 \(A\alpha_{i}\) 等于 A，而后 r 个异于 A。则 \(M^{\prime}/M\) 也同构于诸模 \(A/A\alpha_{n}, \ldots, A/A\alpha_{n-r+1}\) 的直和，其中 \(A\alpha_{n} \subset A\alpha_{n-1} \subset \ldots \subset A\alpha_{n-r+1} \neq A\).

于是 VII 第 16 页命题 2 的推论的条件得到满足：诸理想 \(\mathbf{A}\alpha_{i}\)（\(1 \leqslant i \leqslant n\)）从而由模 \(\mathbf{M}'/\mathbf{M}\) 唯一确定.

由于 L 是 \(M'\) 与 \(L'\) 的直和，故 L/M 是 \(M'/M\) 与 \((L'+M)/M\) 的和，且这个和是直和，因为 \((L'+M) \cap M' = M\)；另一方面，\((L'+M)/M\) 同构于 \(L'/(M \cap L')\)（I，第 41 页，定理 4，c)），即同构于 \(L'\)，这表明 \((L'+M)/M\) 是一个自由模，同构于 \(L/M'\).

推论。--- 主理想整环 A 上的自由模 L 中秩有限的子模 M 在 L 中有补，当且仅当 L/M 是无挠的。

在定理1的记号下，若 L/M 是无挠的，则 \(M = M'\)，且 \(M'\) 在 L 中有补 \(L'\)。反之，若 M 在 L 中有补 \(L'\)，则 L/M 同构于 \(L'\)，它是自由的（VII，第 15 页，推论 2），从而更是无挠的。

注记. --- 可能发生这样的情形：自由模 L 中的无限秩子模 M 可使 L/M 是无挠的，但 M 在 L 中没有补（VII，第 60 页，习题 6，b)）。

定义1.——在定理1的记号与假设下，A 的诸理想 \(A\alpha_{i}\) 称为子模 M 关于模 L 的不变因子。

当 A 是有理整数环 Z 或域 K 上一个未定元的多项式环 K{[}X{]} 时，有一种为 A 的每个理想典范地选取生成元的方式：在 Z 的情形取正整数，在 K{[}X{]} 的情形取首一多项式（VII，第 5 页）。在每一种情形下，不变因子 \(A\alpha_{i}\) 的典范生成元（滥用术语）也称为 M 关于 L 的一个不变因子。

